\section{Bounded-run normalization between primitive phases}
\label{sec:syllable-normalization}

The raw-phase bounds above deliberately excluded normalization. A useful
invariant now includes that boundary for a restricted, explicitly represented
class: each relator is a concatenation of a bounded number of powers of
individual generators. Monomial substitutions preserve this class, and
free and cyclic reduction can be performed by integer addition and endpoint
trimming. The implementation uses this observation at the existing
\texttt{normalize\_relators} move. It retains the general reducer when the
bounded probe cannot establish its precondition.

\subsection{Exact summaries and a conservative domain}

A run list is a sequence $(g_1,e_1),\ldots,(g_s,e_s)$ with positive generator
labels, nonzero signed integer exponents and $g_i\ne g_{i+1}$. It denotes the
freely reduced word $g_1^{e_1}\cdots g_s^{e_s}$. The producer evaluates an
SLP in postorder, storing either an exact list of at most $q$ runs or an
unsupported marker. Terminals have one run. To concatenate two supported
children, it folds the right list into a stack containing the left list:
equal adjacent generator labels cause exponent addition, and zero sums
are removed. The result is stored only if it has at most $q$ runs.

The proof of exactness is induction on the grammar. Each child is freely
reduced; only its concatenation boundary can create cancellation. Repeated
boundary cancellation and merging give the unique freely reduced word.
An unsupported child propagates an unsupported result. In particular, this
is a sufficient condition on \emph{all intermediate words}, not a complete
test for whether the final word has at most $q$ runs. For example,
$(ab)^N((ab)^N)^{-1}$ reduces to the empty word, but its many-run intermediate
is deliberately unsupported. It proceeds to the maintained general reducer.

For a supported root, cyclic reduction trims opposite-sign powers of the
same generator at its two ends. It subtracts the smaller absolute exponent,
removes empty ends, and repeats. Equal-sign endpoint powers are left in
place: the implementation preserves the same anchored word as ordinary
cyclic reduction, rather than choosing a rotated conjugate. This matters
because subsequent proof moves can address word positions.

The production limit is $q=4$. Only roots of expanded length greater than
64 are probed; short roots keep the existing reducer. If there are no long
roots, the probe returns immediately. Before building any new word circuit,
it checks every required long root. A failed probe allocates summaries but
no new SLP word nodes. A successful probe rebuilds long roots from binary
powers and retains all relation slots, including empty and repeated slots.
Short roots require only bounded expanded-word work: with nonempty children
a binary parse of at most 64 letters has at most 127 node occurrences.
The gate does not conceal an exponentially long cancellation problem.

Summaries are cached by immutable arena identifier and run limit. Their
meaning is independent of the current relation slot and live generator set.
Producer and replayer have separate caches and implementations. The replayer
merges a left list against a right deque, then trims endpoints with a deque;
it constructs a left-associated output circuit while the producer balances
the run concatenations. Neither imports the other's arithmetic helper.
The old proof move, version rules and full source reconstruction remain in
force. These routines establish word equality, not an independent knot verdict.

\subsection{Closure under monomial maps}

Suppose the initial relators are given by compact run circuits: each root
concatenates at most $q$ single-generator powers, and each power is encoded
by binary powering. A monomial map sends each generator to an integer power
of one generator, possibly with negative exponent or exponent zero.
It replaces each original run by at most one run. Therefore it cannot
increase the number of runs before free reduction, and free or cyclic
reduction cannot increase that number either.

Every intermediate of the mapped compact circuit also has at most $q$
freely reduced runs. An intermediate inside a power circuit is uniform;
an intermediate concatenation covers at most the original $q$ blocks.
The same is true of the shared powers used to construct forest images:
they are still powers of a single surviving generator. Consequently the
bounded probe succeeds throughout any sequence of these substitutions and
normalizations, if the starting circuit has this form. Its rebuilt output
again has the required compact form. Short roots use only bounded literal
work and also admit compact bounded-size circuits.

Both disjoint primitive projections and unit-coordinate forest quotients
are monomial maps. Initial Wirtinger relators have at most four letters,
so they satisfy the representation condition. This does not imply that
the full search stays in the class: nonmonomial Whitehead changes, overlap
moves and general eliminations can leave it. Nor can an arbitrary grammar
be admitted solely because its \emph{final} normal form happens to have
four runs. The intermediate condition is essential to the implemented
probe and to the stated closure argument.

\subsection{Encoded cost including normalization}

Let $M$ count historical allocated nodes, $S$ the currently reachable nodes,
$m$ the retained root slots, and $B\ge1$ bound represented-length and exponent
bit sizes. Constructing previously absent run summaries takes
$O(SqB)$ bit operations for exponent additions, up to identifier and
dictionary costs. Stored summaries use $O(MqB)$ bits under the same
qualification. Cyclic trimming uses $O(mqB)$ exponent work. Rebuilding
compact roots adds $O(mq(B+1))$ nodes; binary-power construction and
length bookkeeping have polynomial bit cost, not unit-cost arithmetic.
Work counters charge traversal and run processing and poll cancellation;
they do not measure individual big-integer instructions or total process memory.

The compact reachable grammar after a successful normalization satisfies
\[
 S_{j+1}=O\bigl(mq(B_{j+1}+1)\bigr),
\]
where the fixed short-word gate is absorbed in the constant. Historical
nodes are not deleted. A projection or forest batch maps each reachable
input node once and adds $O(r_j(B_j+1))$ image nodes. A following compact
normalization therefore gives the additive allocation bound
\[
 M_{j+1}\le M_j+
 O\bigl(S_j+r_j(B_j+1)+mq(B_{j+1}+1)\bigr).
\]
The distinction between reachable and historical nodes prevents an
unjustified repeated polynomial blow-up. Without an intervening
normalization, the earlier raw-round recurrence still applies.

For disjoint-pair rounds, $B_{j+1}\le 2B_j+O(1)$. Combining raw rounds
with any of these supported normalization boundaries yields an encoded
time bound
\[
 \operatorname{poly}(\text{initial encoded parameters})\,2^{O(d)}
\]
for $d$ such rounds, including polynomial-cost candidate checks and replay.
For mixed forest phases the corresponding bit bound is
$B_{j+1}=O((r_0+1)(B_j+1))$, giving
\[
 \operatorname{poly}(\text{initial encoded parameters})\,
       2^{O(d\log(r_0+1))}.
\]
Here $q$ and the short-word gate are fixed, and the phase count includes
the monomial batches; redundant normalization alone does not change a
normalized state. These extend the earlier raw-phase estimates to this
closed normalized class. They do not bound the number of required phases,
prove that a terminal is reached, or include arbitrary legacy exposure
work. A general quasipolynomial recognition theorem still requires those
missing global ingredients.

\subsection{Correctness audits and measured scope}

Six focused tests cover arbitrary binary parses, changed run limits,
repeated roots, endpoint signs, huge exponents, conservative failure,
monomial closure, independent source replay and resource guards. Conjugates
with exponent $2^{4096}$ normalize with expansion, equality, slicing and
the general reducer disabled. An actual five-braid source with a supplied
powered Whitehead prefix passes complete compressed verification with
all producer normalization helpers disabled. The full maintained suite
passes 998 tests with Regina in 209.031 seconds.

Against the complete pinned package at \texttt{483b7397a108}, eighty
validated diagrams give 240 comparisons across default, projection and
forest modes. All certificates and nondecisions agree; all 159 positive
mode results pass both full source replayers. An independent literal audit
checks all 87,381 signed rank-two words through length eight: 66,580 bounded
probes succeed and 20,801 decline, including the empty-word short gate.
A declined probe makes no negative word or knot claim.

These finite equalities are not a universal statement about search traces.
Equivalent rebuilt words may have different grammar shapes, and later
heuristics or finite resource allowances can depend on allocated nodes
or charged work. Soundness follows from exact normalization and complete
source replay, regardless of those possible search differences.

The complete-recognition run uses nineteen actual diagrams, five measured
rounds and one warm-up per arm. Eight shuffled arms compare the old and
new projection and forest modes, each with an A/A control. All 760
measurements and 152 warm-ups complete; the run takes 88.628 seconds.
The following table retains all inputs. Ratios are medians of paired
old/new times, so values above one favor the new implementation and need
not equal ratios of the separate displayed time medians.

\begin{center}\small
\begin{tabular}{@{}lrrrr@{}}
\toprule
Input & pair old/new ms & ratio & forest old/new ms & ratio \\
\midrule
survivor-00 & 18.586/17.619 & 1.055 & 17.893/18.042 & 0.986 \\
survivor-01 & 15.160/15.053 & 1.004 & 15.120/15.052 & 1.000 \\
survivor-02 & 16.720/16.665 & 1.006 & 16.665/16.619 & 1.005 \\
survivor-03 & 11.587/11.672 & 0.989 & 11.646/11.775 & 0.995 \\
mirror-03 & 25.682/26.593 & 1.004 & 25.636/26.613 & 1.007 \\
survivor-04 & 10.680/10.862 & 0.994 & 10.619/10.822 & 1.010 \\
survivor-05 & 20.959/20.010 & 1.047 & 20.375/20.345 & 0.990 \\
survivor-06 & 8.567/8.150 & 1.059 & 8.331/8.232 & 1.012 \\
survivor-07 & 20.513/21.119 & 0.981 & 20.795/21.110 & 0.994 \\
survivor-08 & 8.470/8.319 & 1.018 & 8.160/8.684 & 0.963 \\
mirror-08 & 21.462/23.429 & 0.955 & 21.038/21.568 & 0.982 \\
survivor-09 & 7.897/7.983 & 0.989 & 7.811/7.859 & 0.994 \\
survivor-10 & 7.000/6.955 & 0.989 & 6.830/7.141 & 0.957 \\
survivor-11 & 8.129/7.896 & 1.030 & 8.577/8.157 & 1.031 \\
gordian & 1595.833/1574.728 & 1.004 & 1546.156/1545.253 & 1.001 \\
trefoil & 0.054/0.054 & 0.987 & 0.055/0.055 & 1.019 \\
figure\_eight & 0.062/0.062 & 1.000 & 0.062/0.062 & 1.002 \\
conway & 1.041/1.019 & 1.024 & 1.055/1.050 & 0.996 \\
kinoshita\_terasaka & 1.045/1.077 & 0.977 & 1.038/1.040 & 0.996 \\
\bottomrule
\end{tabular}
\end{center}


Pair-mode ratios range from 0.955 to 1.059, and forest ratios from 0.957
to 1.031. The old/new pair A/A ranges are 0.983--1.069 and 0.963--1.071;
the old/new forest controls span 0.967--1.054 and 0.919--1.060.
These measurements do not establish an overall recognition speedup.
Gordian has pair/forest ratios 1.004/1.001. In both modes its single
long-root probe declines, visiting 51 new summaries; the source has no
successful bounded normalization. Forest work rises from 1,266,731 to
1,267,256, while its 9,084 nodes and 10,060-byte proof are unchanged.
No other input in this corpus activates a long-root producer probe.

\begin{center}\small
\begin{tabular}{@{}lrrrr@{}}
\toprule
Crossings & pair old/new ms & ratio & forest old/new ms & ratio \\
\midrule
8 & 2.060/1.987 & 1.022 & 1.713/1.732 & 0.984 \\
16 & 5.209/5.266 & 0.988 & 4.151/4.176 & 0.994 \\
32 & 15.305/15.401 & 0.996 & 10.931/11.187 & 0.994 \\
64 & 50.333/50.272 & 0.988 & 33.070/33.878 & 1.002 \\
128 & 186.828/182.397 & 1.026 & 112.551/114.080 & 1.004 \\
\bottomrule
\end{tabular}
\end{center}


The separate checked group stages bypass earlier simplification on actual
stabilized circles. All 200 measurements and 40 warm-ups complete in
10.578 seconds. Every normalization is below the long-word gate; these
rows measure dispatch overhead, not the bounded arithmetic kernel.
The 128-crossing stage has pair/forest paired ratios 1.026/1.004.
Charged forest work nevertheless rises from 170,833 to 178,960 because
the new length preflight is charged on each normalization boundary.
Its 1,017 nodes and complete proof are unchanged. Smaller stages include
regressions. Neither these easy circles nor bypassing earlier simplification
supports a claim about difficult whole-knot recognition.

\begin{center}\small
\begin{tabular}{@{}lrrrrrr@{}}
\toprule
Case & bits & old ms & new ms & ratio & old A/A & new A/A \\
\midrule
conjugate & 16 & 0.368 & 0.301 & 1.331 & 1.020 & 0.749 \\
conjugate & 64 & 1.215 & 0.998 & 1.245 & 0.985 & 1.006 \\
conjugate & 256 & 4.515 & 3.581 & 1.262 & 0.985 & 0.985 \\
conjugate & 1024 & 18.598 & 16.018 & 1.188 & 1.003 & 0.994 \\
conjugate & 4096 & 81.615 & 74.807 & 1.164 & 1.041 & 1.139 \\
source-replay & 16 & 1.391 & 1.476 & 0.944 & 1.033 & 1.019 \\
source-replay & 64 & 4.743 & 5.045 & 0.943 & 1.016 & 1.008 \\
source-replay & 256 & 18.237 & 18.704 & 0.945 & 0.999 & 0.979 \\
source-replay & 1024 & 74.004 & 82.450 & 0.879 & 0.852 & 1.051 \\
source-replay & 4096 & 319.300 & 338.978 & 0.916 & 0.997 & 1.039 \\
wide-cancellation & 64 & 0.568 & 0.774 & 0.731 & 1.020 & 0.985 \\
wide-cancellation & 1024 & 8.590 & 11.585 & 0.742 & 1.002 & 1.010 \\
\bottomrule
\end{tabular}
\end{center}


The final experiment separates twelve controlled cases, with four shuffled
arms and five measured rounds. All 240 measurements and 48 warm-ups
complete within their five-second/20-million-work allowances; the complete
run takes 14.216 seconds. In the table, ``bits'' is $b$ for an exponent
$2^b$. Conjugate rows construct a fresh grammar for
$a^{2^b}b^{2^b+3}a^{-2^b}$ and normalize it. Source-replay rows reconstruct
an actual source PD and verify a complete supplied proof with a powered
prefix; proof discovery is outside the timer. Wide-cancellation rows use
$(ab)^{2^b}((ab)^{2^b})^{-1}$ and exercise the unsupported-intermediate
fallback. Each result is checked against its known word or full source proof.

The conjugate kernel has paired ratios 1.164--1.331. At $b=256$, the
ratio is 1.262 with both A/A controls 0.985. At $b=4096$, old/new work
is 127,040/110,645 and allocated nodes are 16,396/12,295, but its new
A/A ratio 1.139 cautions against attributing the entire 1.164 timing ratio
to the implementation. The smallest case is also noisy, with new A/A 0.749.

Complete supplied-proof replay regresses: ratios range from 0.879 to 0.945.
At $b=4096$, old/new separate medians are 319.300/338.978 milliseconds
and the paired ratio is 0.916. One bounded replay succeeds, but visiting
16,386 summaries adds enough work that the total rises from 438,720 to
463,311; both versions allocate 20,511 nodes. The earlier powered-prefix
operations already populate useful general-reduction caches, so a fresh
normalization pass is not automatically a net saving. At $b=1024$ the
old A/A ratio 0.852 also shows substantial timing variation.

Unsupported wide words regress more clearly, with ratios 0.731 and 0.742.
At $b=1024$, both implementations allocate 2,055 nodes, while charged work
increases from 15,393 to 21,577. The probe adds no word nodes but still
pays to establish that its intermediate summaries are unsupported before
calling the general reducer. These retained negative results limit the
practical claim: the change supplies a closed-class encoded bound and
a faster isolated conjugate kernel, not a universal speed improvement.
Projection and forest recognition remain optional.


The reproduction driver is \path{primitive_power_research/syllables.py}
under \path{../fast}, with \texttt{audit}, \texttt{benchmark},
\texttt{stages} and \texttt{kernels} modes. Each result pins the complete
prior Python package and records current source hashes before and after
execution. Raw shuffled samples, controls, counters, complete certificates
and incomplete outcomes are retained. The supplied-prefix timings measure
verification of a given proof; they do not measure discovering that proof
or recognizing a difficult knot from scratch.
